\chapter{Stage}
\label{cap:stage}

\section{Presentazione del progetto}

\subsection{Il problema}

La revisione manuale di documenti tecnici o legali è un'attività che richiede tempo e attenzione. Contratti, specifiche tecniche e documenti normativi devono rispettare strutture precise, contenere determinate sezioni e conformarsi a \emph{policy}. Affidarsi esclusivamente alla revisione umana espone il processo a errori di disattenzione e rende difficile garantire una standardizzazione costante, soprattutto quando i volumi di documenti da analizzare sono elevati. 

\subsection{La soluzione}

Partendo da questa problematica, è stata sviluppata \textit{AI Document Consistency Platform}: una \emph{web application} progettata per automatizzare il controllo di qualità sui documenti testuali. La piattaforma consente di importare documenti in diversi formati (PDF, Word, Markdown) e di sottoporli a un'analisi su diversi aspetti.

Il sistema adotta un approccio ibrido, composto da due livelli distinti. Il primo livello gestisce i controlli strutturali tramite logiche algoritmiche (\emph{rule-based}): controlla la corretta lunghezza delle sezioni definite nel \emph{ruleset} di configurazione. Il secondo livello delega l'interpretazione semantica del testo ai modelli \gls{llm} Claude di Anthropic, con l'obiettivo di rilevare incoerenze contenutistiche, contraddizioni tra paragrafi, violazioni delle \emph{policy} aziendali e mancanza di sezioni obbligatorie.

Al termine dell'analisi, la piattaforma restituisce un \emph{report} strutturato
che evidenzia le criticità riscontrate e propone le relative correzioni, con l'obiettivo di alleggerire i flussi di revisione e garantire un livello elevato di standardizzazione.

\section{Metodologia di lavoro}
 
Le attività sono state pianificate e monitorate tramite \emph{\gls{jira}}\glsfirstoccur, 
strumento che consente la gestione delle \emph{task} e il tracciamento dei progressi, 
visibile a tutti i membri del gruppo e ai responsabili tecnici.
 
\subsection{Strumenti di supporto}
 
\textbf{Versionamento del codice}: Per il versionamento del codice sorgente è stato adottato \emph{\gls{git}}\glsfirstoccur, con \emph{repository} ospitate su \emph{\gls{github}}\glsfirstoccur. L'utilizzo di \emph{branch} differenti per singola funzionalità ha permesso di mantenere ordine sull'organizzazione del progetto e di gestire il processo di integrazione in modo controllato.\\
 
\textbf{Gestione di progetto}: La gestione delle attività è avvenuta tramite \gls{jira}, attraverso il quale vengono assegnate le \emph{task} ai membri del team, definite le priorità e monitorato l'avanzamento delle funzionalità da implementare. Ogni attività viene stata tracciata con la relativa: descrizione, importanza, stato di completamento e responsabile.\\

\textbf{Comunicazione}: Per la comunicazione interna al team si è utilizzato \emph{\gls{slack}}\glsfirstoccur. Nel contesto dello stage è stato creato un canale dedicato al progetto, all'interno del quale venivano: condivisi aggiornamenti sullo stato di avanzamento, segnalati eventuali problemi ed effettuate richieste di chiarimenti.

\subsection{Tecnologie di supporto}
TODO (AWS, React, NestJS, MongoDB, modelli Claude di AntAnthropic)

\section{Rapporto con l'azienda e con il tutor aziendale}
Lo stage si è svolto interamente in presenza presso la sede di Padova, con orario 
di lavoro dalle 9:00 alle 18:00 dal lunedì al venerdì. \\

Il referente formale dello stage è \textbf{Stefano Dindo}, con il ruolo di \emph{Dev Team Lead}. Il supporto operativo quotidiano è stato invece fornito da \textbf{Michele Massaro}, presente in ufficio quasi ogni giorno, con il quale era possibile confrontarsi sia di persona che tramite \gls{slack} in base alla sua disponibilità.\\

Non sono stati programmati incontri fissi settimanali, ma il rapporto con Michele Massaro è stato continuo: per dubbi o problemi tecnici era possibile ottenere un confronto direttamente durante la giornata, senza dover attendere riunioni formali.

\section{Analisi dei rischi}

Di seguito vengono riportati i principali rischi individuati prima dell'avvio dello sviluppo, con la relativa analisi e le strategie adottate per gestirli.

\rowcolors{2}{tablerowgreen}{white}
\begin{longtable}{L{2.2cm}|L{1.5cm}|L{1.5cm}|L{2.5cm}|L{2.5cm}|L{2.5cm}}
\rowcolor{headergreen}
\thead{Rischio} & 
\thead{Prob.} & 
\thead{Impatto} & 
\thead{Rilevamento} & 
\thead{Mitigazione} & 
\thead{Contingenza} \\
\endfirsthead
\rowcolor{headergreen}
\thead{Rischio} & 
\thead{Prob.} & 
\thead{Impatto} & 
\thead{Rilevamento} & 
\thead{Mitigazione} & 
\thead{Contingenza} \\
\endhead
\endfoot

Tecnologie non conosciute & Alta & Medio & Valutazione delle competenze prima dell'avvio & Studio della documentazione ufficiale & Supporto del tutor aziendale \\
\hline
Ritardo nello sviluppo & Media & Alto & Monitoraggio settimanale tramite \gls{jira} & Pianificazione con priorità definite sin dall'inizio & Eliminazione dei requisiti desiderabili e facoltativi \\
\hline
Comprensione errata dei requisiti & Media & Alto & Revisione con il responsabile prima dello sviluppo & Analisi con il responsabile durante la raccolta dei requisiti & Correzione dei requisiti con il tutor prima dell'implementazione \\
\hline
Eterogeneità dei formati documentali & Media & Medio & Verifica durante la fase di sviluppo dei \emph{parser} & Utilizzo di librerie consolidate per ogni formato (PDF, Word, Markdown) & Limitare i formati supportati ai più diffusi \\
\hline
Indisponibilità o costo elevato delle \gls{api} Anthropic & Bassa & Alto & Monitoraggio dell'utilizzo durante la fase di test & Limitare le chiamate alle \gls{api} durante lo sviluppo & Utilizzo di un piano di \emph{fallback} con risposte semplificate \\
\rowcolor{white}\caption{Analisi dei rischi}
\label{tab:rischi}
\end{longtable}